\chapter{Практика}

Построим модель линейной регрессии. За целевую переменную возьмем параметр PRP.

\section{Матрица корреляции}

\begin{figure}[H]
    \centering
    \includegraphics{corr.png}
    \caption{Матрица корреляции}
\end{figure}

Матрица корреляции показывает зависимости между признаками. По матрице выше можно сказать, что почти все признаки неплохо коррелируют друг с другом.

\section{Сравнение реальных значений и предсказанных}

\begin{figure}[H]
    \centering
    \includegraphics{main.png}
    \caption{Сравнение реальных значений и предсказанных}
\end{figure}

Красной линией показано идеальное расположение ответов модели, когда предсказанные значение точно совпадают с рельными. \\

Здесь есть небольшие расхождения ответов с реальными данными. В начале расхождения меньше, чем в конце. Скорее всего это связано с выбросами

\begin{figure}[H]
    \centering
    \includegraphics{err.png}
    \caption{Ошибка}
\end{figure}

Ошибку лучше видно на графике выше. В начале ошибка намного меньше.

\section{Распределение ошибки}

\begin{figure}[H]
    \centering
    \includegraphics{distrb.png}
    \caption{Распределение ошибки}
\end{figure}

Ошибка имеет около-нормальное распределение.

\endinput